\documentclass{article}

\usepackage[spanish]{babel}
\usepackage{amsmath,amssymb,amsthm}

\theoremstyle{remark}
\newtheorem{remark}{Observación}

\title{Historia de la Arquitectura de Computadores}
\author{Benjamín Rivas Beltrán}
\date{\today}


\begin{document}

\maketitle

\section{Computación Pre-mecánica}

\begin{itemize}
    \item De contar con los dedos.
    \item A contar con piedrecillas.
    \item A marcar líneas en muros.
    \item A marcar líneas en huesos.
    \item A marcar líneas en la arena.
\end{itemize}

\textbf{Pregunta interesante:} ¿Hay alguna otra especie a diferencia del ser humano que puede contar?

\section{Computación Mecánica}

\begin{itemize}
    \item Abaco.
    \item Los huesos de Napier y el logaritmo.
    \item La regla de cálculo de Oughtred.
    \item Reloj calculador de Schickard.
    \item La Pascalina.
    \item La calculadora paso a paso de Leibniz.
    \item Telares controlados por tarjetas perforadas de Jacquard.
    \item Charles Babbage, el padre de la computadora.
    \item Ada Lovelace, la primera programadora de la historia.
    \item La máquina tabuladora de Hollerith.
    \item La máquina de Turing.
    \item John Vincent Atanasoff y Clifford Berry, creadores de la primera computadora electrónica digital.
\end{itemize}

\begin{remark}
Alan Turing, al conseguir decodificar los mensajes de la máquina Enigma, sentó las bases de lo que hoy conocemos como la \textit{Ingeniería inversa}.
\end{remark}

\begin{remark}
John Vincent Atanasoff demostró que los circuitos electrónicos son isomorfos a los circuitos mecánicos, lo que dio paso a los dispositivos electromecánicos. 
\end{remark}

\end{document}